\subsection*{Preliminary lemma}

% Nur kurze zusammenfassung der Notation zur Übersicht
Notation:
\begin{itemize}
	\item $\curve$: parametric curve, $\len$: length of curve, $\sublen$: length of subcurve from $a$ to $x$
	\item $\metric_{\infty}$: $\infty$-metric, if no subscript then $\metric$ is the additive difference metric
	\item $\metric$: metric, $\metr{a}{b}$ metric distance between $a$ and $b$, $\metro{a}:= \metr{0}{a}$
	\item $\ball{r} := \{x \in \mathbb{R}^{n}: \metro{x} < r\} $: the open Ball with radius $r$ around $0$ wrt. the metric $\metric$,
	\item $\cball{r} := \{x \in \mathbb{R}^{n}: \metro{x} \leq r\} $: the closed Ball with radius $r$ around $0$ wrt. the metric $\delta$,
	\item $\|x,y\|_\infty := \max_{i \leq n} \{|x_i - y_i|\}$: the maximum metric. We just use $\infty$ to shorten the notation, e.g. in $\ball[\infty]{r}$.
	\item $\co(A)$: convex Hull of $A$.
	\item $\extr(A)$: extreme Points of $A$.
\end{itemize}

In this section we will prove the main lemma, which is needed for the proof of the main result \cref{lp-uniqueness-theorem}.
First, we will state the lemma, then give some remarks on the lemma and prove one statement needed for the lemma and then finally prove it.
The proof is divided into 5 parts and a whole subsection which is devoted to showing a statement that is needed in the proof.
Everything in this section will be assumed as in this Theorem. For example, if we write $\metric$, this implies not any metric but the metric from the lemma with its properties.
Moreover all topological terms(open, closed, continuous) are meant with respect to the natural topology of $\mathbb{R}^{n}$ (i.e. the topology induced by the $\metric_{\infty}$ and the euclidean metric), not the topology induced by $\metric$ or the induced topology on $S$. This will be especially important for the subsection on convexity of the balls.
We will begin with the statement:

% TODO: change to thm and edit text everywhere, find fitting name for this theorem
\begin{lem}[Main Lemma] \label{main-lemma}
	Let $S$ be an open convex subset of $n$-dimensional Euclidean space. Let $\metric$ be a metric on $S$ that satisfies intradimensional subtractivity \footnote{i.e. $\metric$ takes the form $\metr{x}{y} = F(|x_1 - y_1|, \dots, |x_n - y_n|)$.} with respect to a continuous function $F: \mathbb{R}_{\geq 0}^n \rightarrow \mathbb{R}_{\geq 0}$ and linear coordinates \footnote{Linear coordinates means that the vector representations are cartesian as opposed to, for example, polar, where each coordinate entry is not the value on the corresponding factor of the cartesian product.} in $S$. If $\langle S, \metric \rangle$ is a metric with additive segments, then $\metric$ is homogeneous, i.e., for any $x, z \in S$ and any real $0 \leq t \leq 1$,
	\begin{align*}
		\metric[x, x + t(z - x)] = t\metric(x, z).
	\end{align*}

\end{lem}

Now we will continue with a remark that makes it possible to move the set such that $0 \in S$ by using the translation invariance of $\metric$.

\begin{rem}[Translation invariance and domain of $\metric$ and $F$] \label{lemma-rem}

	\begin{enumerate}[(i)]
		\item $\metric$ is translation-invariant, i.e. $\metr{x+y}{z+y} = \metr{x}{z}$: \label{lemma-trans-invariance}
		      \begin{align*}
			      \metr{x+y}{z+y} & = F[|x_1 + y_1 - (z_1 + y_1)|, \dots, x_n + |y_n - (z_n + y_n)|] \\
			                      & = F[|x_1 - z_1|, \dots, |x_n - z_n|]                             \\
			                      & = \metr{x}{z}
		      \end{align*}
		\item Because $\metric$ is invariant to translation in its domain, we can move $S$ and redefine the metric on the shifted set. We do this by defining $S' = \{s - x : s \in S\}$ for $x \in S$ and $\metric'(a,b) := \metr{a+x}{b+x}$ for $a, b \in S'$. Since $0 \in S'$, we also can introduce the following shorthand Notation $\metric'(x') := \metric'(0, x')$ for $x' \in S'$.
		      From now on, we will assume that $0 \in S$, because by the argument just provided, we can shift $S$ and $\metric$ such that $0 \in S$.
		      Thus we can proceed to proof the statement using the shorthand notation, by which we need to show that:
		      $$ \metro{t \cdot  z} = t \cdot \metro{z} \qquad \forall t \in [0, 1] \text{ and } \forall z \in S' .$$ \label{rem-lemma-translate-to-zero}
          %TODO: Bei (iii) auf S.5 fehlt mir auch ein bisschen das Argument. Irgendeine Voraussetzung an F muss da noch eingehen. Nur weil \delta eine Metrik ist, muss ja nicht der Definitionsbereich ein kartesisches Produkt sein, oder uebersehe ich da etwas? 
		\item Since $\metric$ is defined by $\metr{x}{y} = F(|x_1 - y_1|, \dots, |x_n - y_n|)$ and $\metric$ is a metric, $F(x) = 0$ iff $x=0$ must hold. Moreover the function $F$ has to be defined on a set, where at least the "greatest possible difference" in coordinate $i$ is a possible $i$-th argument. This means $F$ is defined on a cartesian product of half-open, real intervals of the form $[0,a_i)$ where $$a_i = \sup \left \{ |x_i - y_i| : (x_1, \dots, x_n), (y_1, \dots, y_n) \in S \right \} \in \mathbb{R}_{> 0} \cup \{\infty\}.$$ These $a_i$ are these "greatest possible differences" on the $i$-th coordinate.\label{rem-lemma-domains}
		\item  Moreover, due to the continuity of $F$ and $| \cdot |$, the metric $\metric$ is continuous. \label{rem-lemma-continuity}
	\end{enumerate}
\end{rem}


Before moving to the proof, we will fist show two statements that we will need right at the start.
The first one asserts that the distance between a compact and a closed set is greater zero.

\begin{lem} \label{lemma-distance}
	If $I \subset \mathbb{R}^{n}$ is compact, $Y \subset \mathbb{R}^{n}$ is closed and $I$ and $Y$ are disjoint, then there exists $\varepsilon > 0$ such that $\metr[\infty]{p}{y} > \varepsilon$ for all $p \in I$ and $y \in Y$.
\end{lem}

\begin{proof}
	We will prove this statement by contradiction.
	Assume that for all $\varepsilon > 0$ there exists $p \in I$ and $y \in Y$ such that $\metr[\infty]{p}{y} < \varepsilon$. This means, we can construct two sequences $p_n \in I$ and $y_n \in Y$ such that $\metr[\infty]{p_n}{y_n} \to 0$ as $n \to \infty$.
	$I$ is compact, thus there exists a convergent subsequence $p_{n_m}$ such that $p_{n_m} \to p \in I$ for $n \rightarrow \infty$.
	With the triangle inequality we get
	$$\metr[\infty]{p}{y_{n_m}} \leq \metr[\infty]{p}{p_{n_m}} + \metr[\infty]{p_{n_m}}{y_{n_m}} \to 0 \text{ as } m \to \infty.$$
	Hence,  $p$ is a limit point of $Y$, and since  $Y$  is closed,  $p \in Y$ must hold which contradicts the assumption.
\end{proof}

The second one allows us to make arbitrary small $\metric$-balls which we will combine with the statement before to assure a distance in $\metric$ and not $\metric_{\infty}$.

\begin{lem}[arbitrary small $\metric$-Balls] \label{lemma-small-delta-balls}
	Let $\varepsilon > 0$ such that $\ball[\infty]{\varepsilon} \subset S$. Then there exists $\alpha > 0$ such that $\ball{\alpha} \subset \ball[\infty]{\varepsilon}$.
\end{lem}

\begin{proof}
	Since $\metric$ is a metric defined by $\metric(x,y) = F[|x_1 - y_1|, \dots, |x_n - y_n|]$, we can define the $j$-th partial function of $F$ by setting all other arguments to zero, i.e. $F^j(u_j) := F(0, \dots, u_j, \dots, 0)$.

	$F$ is strictly increasing in every argument implies that $F^j$ is increasing and the continuity of $F$ directly translates to $F^j$ being continuous.
	Moreover, due to \crefpart{lemma-rem}{rem-lemma-domains}, $F^j(u_j) = 0$ if and only if $u_j = 0$ holds.

	Together, this means that $F^j$ is an increasing, continuous real function with $F^j(0) = 0$.


	By setting $\alpha_j = F^j(\varepsilon)$ we get by monotonicity, that $F^j(u_j) < \alpha_j \implies u_j < \varepsilon$. Note, that $F^j(\varepsilon)$ is defined, because  $\ball[\infty]{\varepsilon} \subset S'$ implies that $\varepsilon$ is a possible distance for each coordinate. (See also \crefpart{lemma-rem}{rem-lemma-domains}. Using the $a_i$ from the remark, we get that $\varepsilon < a_i$ holds for all $i$ and thus $F^j(\varepsilon)$ is defined for all $j = 1, \dots n$.)

	Doing this for $j = 1, \dots, n$ yields $\alpha_1, \dots, \alpha_n$.
	By setting $\alpha = \min \{\alpha_1, \dots, \alpha_n\}$ the above holds for every partial function $F^j$ which yields the boundary for the $\infty$ Metric:
	$$
		F^j(u_j) < \alpha \leq \alpha_j \implies u_j < \varepsilon \quad \forall j = 1, \dots, n \implies \|u\|_{\infty} < \varepsilon.
	$$

	Using the monotonicity of $F$ we get
	$$
		\metro{u} = F(u) < \alpha \implies F^j(u_j) < \alpha \implies \|u\|_\infty < \varepsilon.
	$$
\end{proof}


Now, we will move to the proof of the main lemma. It is suggested to first read the proof to get the main ideas and a big picture understanding, and then read the next section which covers the convexity of spheres which will be of great importance.
First we will restate the lemma in the form, in which we will prove it.

\begin{lem*}[Restatement of the Lemma]
	Let $S$ be an open convex subset of $n$-dimensional Euclidean space with $0 \in S$. Let $\metric$ be a metric on $S$ that satisfies intradimensional subtractivity with respect to linear coordinates in $S$. If $\langle S, \metric \rangle$ is a metric with additive segments, then $\metric$ is homogeneous, i.e., for any $z \in S$ and any real $0 \leq t \leq 1$,
	\begin{align*}
		\metro{tz} = t\metro{z}.
	\end{align*}
\end{lem*}

The proof consists of 4 parts.
First we will find a radius such that all balls with this radius on line from $0$ to $z$ are in $S$.
Then we will prove this lemma for $t = \frac{1}{m}$ for $m \in \mathbb{N}$ by proving the two inequalities.
After this, we will expand this statement to the rationals and then use continuity to expand it to all real numbers between $0$ and $1$.

\begin{proof}

	\setcounter{proofpart}{-1}


	%-------------- Proof part 0 ----------------------------%
	\begin{proofpart}[Asserting a common distance for all points on the line segment] \label{lemma-part0}
		Let $z \in S$.
		Since $S$ is open, $\mathbb{R}^{n} \setminus S$ is closed. The set $I := \{t z : t \in [0,1]\}$ is compact, since continuous images($t \mapsto t z$) of compact sets ($t \in [0,1]$) are compact\todo{ref?}.
		Thus, by \cref{lemma-distance} there exists a distance $\varepsilon > 0$ such that $\metr[\infty]{p}{y} < \varepsilon$ for all $p \in I$ and $y \in \mathbb{R}^{n} \setminus S$.
		This means, for all points on the line from $0$ to $z$ the ball with radius $\varepsilon$ is a subset of $S$, i.e. for all $t \in [0,1]: \ballp[\infty]{\varepsilon}{tz} \subset S$.
    By \cref{lemma-small-delta-balls}, there exists $\alpha > 0$ such that $\ball{\alpha} \subset \ball[\infty]{\alpha}$, and by translation invariance this holds for all points on the line: $\ballp{\alpha}{tz} \subset S$ for all $t \in [0,1]$.
		Now let $m \in \mathbb{N}$ such that $\frac{1}{m} \metro{z} < \alpha$.
		In particular, this means that $\ball{\alpha}$ and therefore, $\ball{\frac{1}{m} \metro{z}}$ is bounded (with respect to the $\infty$ metric).

	\end{proofpart}



	%-------------- Proof part 1 ----------------------------%
	\begin{proofpart}[$t \metro{z} \leq \metro{t z}$ for $t = \frac{1}{m} < d$ with $m \in \mathbb{N}$]\label{lemma-part1}
		Let $\z[k] := \frac{k}{m} \cdot z$ for $k =0, \dots, m$. $\z[k] \in S$ for all $k$ due to convexity of $S$.

		By applying translation invariance, we can calculate the distance between two consecutive points:
		$$
			\metr{\z[k-1]}{\z[k]} = \metr{\frac{k-1}{m} z}{\frac{k-1}{m} z + \frac1m z} = \metro{\frac{1}{m} z} \quad \forall k = 1, \dots, m.
		$$

		Applying the triangle inequality $m-1$-times and using the above result yields

		$$
			\metro{z} \leq \sum_{k=0}^{m-1} \metr{\z[k]}{\z[k+1]} = m \cdot \metro{\frac{1}{m} z}.
		$$

		Finally, dividing the right side by $m$ yields $\frac{1}{m} \metro{z} \leq \metro{\frac{1}{m} z}$, which concludes the first part.
	\end{proofpart}


	%-------------- Proof part 2 ----------------------------%
	\begin{proofpart}[$t \metro{z} \geq \metro{t z}$ for $t = \frac{1}{m} < d$ with $m \in \mathbb{N}$] \label{lemma-part2}
		Since $\metric$ is a metric with additive segments, there exists a segment from $0$ to $z$ with length $\metro{z}$. \cref{sublen-cont} yields, that the arc length is a continuous, non-decreasing function. This means, we can divide the segment into $m$ subsegments with length $\frac{1}{m} \metro{z}$ and choose $\y[0], \dots, \y[m]$ equally spaced on the segment such that $\metr{\y[i]}{\y[i-1]} = \frac{1}{m} \metro{z}$ and $\y[0] = 0$ and $\y[m] = z$.
		Using this, we can define $\x := \y[i] - \y[i-1]$. Translation invariance yields $\metro{\x} =\metr{\y[i]}{\y[i-1]} = \frac{1}{m} \metro{z}$ and, due to the choice of $m$ in \cref{lemma-part0} $\x \in S$ holds. Note, that this is still true, if we increase $m$.

		Further, we can observe that $\frac{1}{m} z$ is a convex combination of $\x$:
		$$
			\sum_{i=1}^m \frac{1}{m} \x = \sum_{i=1}^m \frac{1}{m} (\y[1] - \y[0] + \y[2] - \y[1] + \dots + \y[m]) = \frac{1}{m} (-\y[0] + \y[m]) = \frac{1}{m} z
		$$

		We showed before, that $\metro{\x} = \frac{1}{m}\metro{z}$, which implies $\x \in  \cball{\frac{1}{m}\metro{z}}$.
		By \cref{lemma-part0} $\ball{\frac{1}{m} \metro{z}}$ is bounded. Therefore we can apply \cref{lem-metr-spheres-convex} from which we get that spheres are convex. This means, that the spheres are closed under convex combination, which implies that $\frac{1}{m} z \in \cball{\frac{1}{m} \metro{z}}$ must hold. Thus $\metro{\frac1m z} \leq \frac1m \metro{z}$ which concludes the second part of the proof.

	\end{proofpart}



	Now, by combining \cref{lemma-part1} and \cref{lemma-part2}, we have established the equality for $t = \frac{1}{m}$. Note, that the equality also holds for all $\mathbb{N} \ni m' > m$ since we only needed $m$ to be big enough in \cref{lemma-part2} for $\x$ to be in $S$ and for convexity of the spheres. Thus, by making $m$ bigger, the crucial points in \cref{lemma-part2} and therefore the equality are still true.
	Next, we will proceed to extending the equality from $t = \frac{1}{m}$ to $t =\frac{k}{m}$ for $k = 0, \dots, m$.



	%-------------- Proof part 3 ----------------------------%
	\begin{proofpart}[$t \metro{z} = \metro{t z}$ for $t = \frac{k}{m}$ and $t \in \{0, \dots, m\}$] \label{lemma-part3}
		% To show this part, we will first use translation invariance to prove $\metro{\frac1m z} = \frac1m \metro{z}$ from all points $\z$.
		For $\z[1]$ we already have $\metr{0}{\z[1]} = \frac{1}{m} \metr{0}{z}$.
		Now, for consecutive points, by translation invariance, it holds that
		$$
			\metr{\z[k-1]}{\z[k]} = \metr{\z[k-2]}{\z[k-1]} = \dots = \metr{\z[1]}{\z[0]} = \metr{\z[1]}{0} = \frac{1}{m} \metr{0}{z}.
		$$

		Using this and the triangle inequality, we get
		\begin{align*}
			\metr{0}{z}          & \leq \metr{0}{\z[1]} + \metr{\z[1]}{\z[2]} + \dots + \metr{\z[m-1]}{\z[m]} \\
			                     & = n \cdot \frac{1}{n} \metr{0}{z}                                          \\
			                     & = \metr{0}{z}                                                              \\[1ex]
			\implies \metr{0}{z} & = \metr{0}{\z[1]} + \metr{\z[1]}{\z[2]} + \dots + \metr{\z[m-1]}{\z[m]}
		\end{align*}\\

		For the distance $\metr{0}{\z[2]}$, we get the following two inequalities:
		\begin{enumerate}
			\item $$ \metr{0}{\z[2]} \leq \metr{0}{\z[1]} + \metr{\z[1]}{\z[2]} = \frac{2}{m} \metr{0}{z} $$
			\item \begin{align*}
				      \metr{0}{z}     & \leq \metr{0}{\z[2]} + \metr{\z[2]}{z}                                                  \\
				      \metr{0}{\z[2]} & \geq \metr{0}{z} - \underbrace{\metr{\z[2]}{z}}_{\leq \sum_i \metr{\z[2+i]}{\z[2+i+1]}} \\
				                      & \geq \metr{0}{z} - \sum_{i=0}^{m-3} \metr{\z[2+i]}{\z[2+i+1]}                           \\
				                      & = \metr{0}{z} - \frac{m-2}{m} \metr{0}{\z}                                              \\
				                      & = \frac{2}{m} \metr{0}{z}
			      \end{align*}
		\end{enumerate}

		These two inequalities yield $\metr{0}{\z[2]} = \frac{2}{m} \metr{0}{z}$.
		By repeating this for $\metr{0}{\z}$ for every $k = 2, \dots, m-1$, we get the equality for every rational number for $m$ big enough. However every rational number with a denominator smaller than $m$ can be represented by another rational number with denominator greater than $m$. This yields, that the result is true for every rational $t$.
	\end{proofpart}




	%-------------- Proof part 4 ----------------------------%
	\begin{proofpart}[$t \metro{z} = \metro{t z}$ for $t \in {[0, 1]}$]
		The final step can be completed using continuity of the metric. Now let $t \in [0,1]$ and $t_n \in \mathbb{Q}$ with $t_n \rightarrow t$ for $n \rightarrow \infty$.
		Then
		\begin{align*}
			\metro{t z} & = \metro{ \lim_{n \rightarrow \infty } t_n z }                                                 \\
			            & = \lim_{n \rightarrow \infty } \metro{t_n z} \tag{continuity}                                  \\
			            & = \lim_{n \rightarrow \infty } t_n \metro{z} \tag{$t_n \in \mathbb{Q}$ and \cref{lemma-part3}} \\
			            & = t \metro{z}
		\end{align*}
	\end{proofpart}
\end{proof}





